% ============================================================================
% 3. DW-BENCH DESIGN
% ============================================================================
\section{DW-Bench: Benchmark Design}

\subsection{Datasets}

% TABLE 1: Dataset Overview
\begin{table}[t]
\centering
\caption{DW-Bench dataset statistics. ``Silos'' = disconnected
components in the combined FK+lineage graph.}
\label{tab:datasets}
\small
\begin{tabular}{@{}lcccccc@{}}
\toprule
Dataset & Domain & Tbl & FK & Lin. & Qs & Silos \\
\midrule
AdventureWorks & Retail   & 102 & 136 & 39 & 208 & 11 \\
TPC-DS         & Analytics&  24 &  70 &  0 & 127 &  1 \\
TPC-DI         & ETL      &  35 &  29 & 21 & 181 &  2 \\
OMOP CDM       & Health   &  37 &  74 & 21 & 158 &  3 \\
Syn-Logistics  & Supply   &  64 &  96 & 35 & 372 &  5 \\
\midrule
\textbf{Total} &          & \textbf{262} & \textbf{405} & \textbf{116} & \textbf{1046} & --- \\
\bottomrule
\end{tabular}
\end{table}

We select five datasets shown in Table~\ref{tab:datasets} representing diverse data warehouse topologies:

\textbf{AdventureWorks}~is a Microsoft reference data warehouse with 102
tables spanning OLTP and DW layers, connected by 136 FK edges and 39
lineage (derived\_from) edges. Its dual-layer structure makes it ideal for
lineage impact questions.

\textbf{TPC-DS}~is the industry standard analytics benchmark with a star
schema of 24 tables and 70 FK edges. Its single connected component and
absence of lineage edges tests pure FK-based reasoning.

\textbf{TPC-DI}~is the data integration benchmark with 35 tables modeling
an ETL pipeline from staging to warehouse, with 21 lineage edges
representing the transformation flow.

\textbf{OMOP CDM}~(Observational Medical Outcomes Partnership Common Data
Model)~\cite{ohdsi2019omop} is a healthcare standard with 37 tables, 74
FK edges, and 21 lineage edges across 3 connected components (clinical
data, vocabulary tables, and metadata).

\textbf{Syn-Logistics}~is our own synthetic supply-chain schema: 64
tables, 5 connected components (carrier, procurement, healthcare, finance,
HR).  We built it to fix a statistical-power problem in the real-world
corpus: every subtype gets $n \geq 20$ questions.  Table names come from
a domain dictionary so obfuscation does not collapse semantics.

\subsection{Schema Graph Representation}

Each dataset is represented as a heterogeneous graph $G = (V, E_{fk} \cup
E_{lin})$, where $V$ denotes the set of table nodes, $E_{fk}$ denotes the
set of foreign-key edges, and $E_{lin}$ denotes the set of data-lineage
edges, using PyTorch Geometric~\cite{fey2019pyg}
\texttt{HeteroData}. Each node represents a table and is associated with six structural features: in-degree, out-degree, normalized
degree, lineage degree, betweenness centrality, and PageRank. Edges are
typed as either \texttt{fk\_to} (foreign key) or \texttt{derived\_from}
(data lineage). Lineage edges form a strict DAG (directed acyclic graph):
each \texttt{derived\_from} edge points from a downstream DW table to its
upstream source, with no cycles or self-references. All traversal
algorithms respect edge typing: FK queries use only \texttt{fk\_to}
edges, lineage queries use only \texttt{derived\_from}, and
\texttt{combined\_impact} queries compose both types sequentially.
FK edges are treated as \textbf{directed} (parent$\to$child) for
path and hop-count queries (\texttt{join\_path}, \texttt{hop\_count},
\texttt{direct\_fk}) and as \textbf{undirected} for connectivity
queries (\texttt{membership}, \texttt{connected}, \texttt{isolation},
\texttt{count}, \texttt{full\_enum}).Lineage edges are always directed (downstream$\to$upstream),
following the \texttt{derived\_from} convention defined above.
Foreign keys are treated as bidirectional (undirected) explicitly for
connectivity and silo-detection subtypes, as an FK implies an inherent
structural relationship that can be traversed logically in either direction
(e.g., from a parent lookup table to child fact rows, or vice versa).

For path-based questions, we validate alternative valid shortest paths
against the FK adjacency matrix, ensuring EM is robust to tie-breaking.

\subsection{Question Generation}
\label{sec:qa-generation}

Questions are generated deterministically from the graph structure using
standard graph algorithms (BFS, connected component detection)
implemented with NetworkX~\cite{networkx2008}, ensuring ground-truth
answers are provably correct. We define three categories with 13 subtypes as shown in Table~\ref{tab:taxonomy}. 

\paragraph{Difficulty assignment.}
Difficulty labels are assigned a priori based on structural complexity,
following the precedent of Spider~\cite{yu2018spider} (which uses SQL AST
depth) and BIRD~\cite{li2023bird}.  Easy subtypes require a single graph
lookup (one hop, one edge type).  Medium subtypes require multi-hop
traversal within a single edge type.  Hard subtypes require either
multi-hop transitive closure or composition across both FK and lineage
edge types.  These labels are fixed properties of the question structure,
independent of any model's empirical performance.

% TABLE 2: Question Taxonomy
\begin{table}[t]
\centering
\caption{Question taxonomy. Difficulty is assigned based on the number of
reasoning hops and edge types required.}
\label{tab:taxonomy}
\small
\begin{tabular}{@{}llclr@{}}
\toprule
Category & Subtype & Diff. & Description & \#Qs \\
\midrule
\multirow{5}{*}{\rotatebox{90}{\small Lin.}} 
  & forward         & Easy & Direct lineage targets   & 72 \\
  & reverse         & Easy & Source tables for DW table & 71 \\
  & transitive      & Hard & Multi-hop lineage chains  & 33 \\
  & combined\_impact & Hard & Lineage + FK dependents  & 69 \\
  & multi\_source   & Med. & Tables with 3+ sources   & 49 \\
\midrule
\multirow{3}{*}{\rotatebox{90}{\small Rt.}}
  & direct\_fk      & Easy & FK adjacency check       & 100 \\
  & join\_path      & Med. & FK shortest path         & 342 \\
  & hop\_count      & Med. & Path length              & 124 \\
\midrule
\multirow{5}{*}{\rotatebox{90}{\small Silo}}
  & count           & Easy & Number of components     & 30 \\
  & membership      & Hard & Which component has X?   & 38 \\
  & isolation       & Med. & Is table X isolated?     & 60 \\
  & connected       & Med. & Are X and Y connected?   & 29 \\
  & full\_enum.     & Hard & List all tables in silo  & 29 \\
\bottomrule
\end{tabular}
\end{table}

\subsection{Obfuscation Protocol}

To distinguish genuine topology-based reasoning from reliance on
surface lexical cues and schema-name memorization,
we generate an \emph{obfuscated} variant for each dataset.  Table names
are replaced with random identifiers (\texttt{Table\_A},
\texttt{Table\_B}, etc.) using a deterministic mapping.
This prevents models from exploiting semantic signals in
names such as \texttt{Customer} or \texttt{Invoice} to infer
relationships without genuine structural reasoning.
Questions and answers are updated accordingly with
word-boundary-aware replacement to avoid corrupting natural
language phrases (e.g., the word ``relationship'' in
``foreign key relationship'' is preserved).
